\chapter{Bilangan Negatif}\label{ch:g7:negatives}

Suhu di bawah \omterm{def:g1:counting:zero}{nol}, lantai di bawah tanah, uang yang terutang: banyak
besaran turun di bawah $0$. Bilangan negatif memperluas garis bilangan
ke kiri dari \omterm{def:g1:counting:zero}{nol}. Bab ini mengajarkan cara membandingkannya serta cara
menjumlahkan dan mengurangkannya --- perkaliannya menunggu
\cref{ch:g8:negprod}.

\section{Bilangan bertanda}

\begin{definition}[Bilangan bertanda]\label{def:g7:negatives:def}
\emph{Bilangan bertanda}\index{bilangan bertanda} tersusun dari sebuah
tanda ($+$ atau $-$) dan sebuah \emph{jarak ke \omterm{def:g1:counting:zero}{nol}}: $+3$ dan $-3$
sama-sama berjarak $3$ dari $0$, di sisi yang berlawanan. Bilangan
positif biasanya ditulis tanpa tandanya: $+3 = 3$. Bilangan $+3$ dan
$-3$ saling menjadi \emph{lawan}\index{lawan}.
\end{definition}

\begin{omfigure}
\begin{tikzpicture}[scale=1.15]
  \draw[-Stealth] (-5.6,0) -- (5.6,0);
  \foreach \x in {-5,...,5} \draw (\x,-0.08) -- (\x,0.08);
  \foreach \x in {-5,...,5} \node[below, font=\small] at (\x,-0.1) {$\x$};
  \fill[omThm] (-3,0) circle (2pt);
  \fill[omDef] (3,0) circle (2pt);
  \draw[Stealth-Stealth, omExo] (-3,0.35) -- (0,0.35)
    node[midway, above, font=\small] {$3$};
  \draw[Stealth-Stealth, omExo] (0,0.35) -- (3,0.35)
    node[midway, above, font=\small] {$3$};
\end{tikzpicture}

{\small Garis bilangan yang diperpanjang ke kiri: $-3$ dan $+3$ adalah
bilangan yang saling berlawanan, berjarak sama dari $0$.}
\end{omfigure}

\begin{method}[Membandingkan bilangan bertanda]\label{met:g7:negatives:compare}
Pada garis bilangan, makin kecil berarti makin ke kiri. Dalam praktik:
\begin{enumerate}
  \item bilangan negatif selalu lebih kecil daripada bilangan positif:
        $-7 < 2$;
  \item di antara dua bilangan positif, berlaku urutan yang biasa;
  \item di antara dua bilangan negatif, yang jaraknya ke \omterm{def:g1:counting:zero}{nol}
        \emph{lebih besar} justru \emph{lebih kecil}: $-7 < -2$.
\end{enumerate}
\end{method}

\begin{example}\label{ex:g7:negatives:compare}
Urutkan dari yang terkecil ke yang terbesar: $3$; $-5$; $0$; $-1.5$;
$2.8$. Yang negatif lebih dulu, yang paling negatif di depan:
\[
-5 < -1.5 < 0 < 2.8 < 3 .
\]
\end{example}

\section{Menjumlahkan bilangan bertanda}

Bayangkan \omterm{def:g7:negatives:def}{bilangan bertanda} sebagai perpindahan sepanjang garis: $+5$
berarti ``berjalan $5$ langkah ke kanan'', $-3$ berarti ``berjalan $3$
langkah ke kiri''. Menjumlahkan bilangan berarti merangkai jalannya.

\begin{method}[Menjumlahkan dua bilangan bertanda]\label{met:g7:negatives:add}
\begin{enumerate}
  \item \emph{Tandanya sama}: jumlahkan jaraknya, pertahankan tanda
        bersamanya: $(-3) + (-4) = -7$.
  \item \emph{Tandanya berlawanan}: kurangkan jarak yang lebih kecil
        dari yang lebih besar, lalu ambil tanda yang lebih besar:
        $(-7) + (+2) = -5$ dan $(+7) + (-2) = +5$.
  \item \omterm{def:g7:negatives:def}{Lawan} saling meniadakan: $(+4) + (-4) = 0$.
\end{enumerate}
\end{method}

\begin{omfigure}
\begin{tikzpicture}[scale=1.05]
  \draw[-Stealth] (-5.4,0) -- (4.4,0);
  \foreach \x in {-5,...,4} \draw (\x,-0.07) -- (\x,0.07);
  \foreach \x in {-5,-3,-1,0,2,4}
    \node[below, font=\small] at (\x,-0.09) {$\x$};
  \draw[-Stealth, omThm, very thick] (0,0.35) -- (-3.5,0.35);
  \draw[-Stealth, omDef, very thick] (-3.5,0.75) -- (-1.5,0.75);
  \node[omThm, font=\small] at (-1.75,0.18) {$-3.5$};
  \node[omDef, font=\small] at (-2.5,1.0) {$+2$};
  \fill (-1.5,0) circle (2pt) node[below=6pt, font=\small] {$-1.5$};
\end{tikzpicture}

{\small Menjumlahkan sebagai berjalan: mulai dari $0$, jalan $-3.5$
lalu jalan $+2$ mendarat di $(-3.5) + 2 = -1.5$.}
\end{omfigure}

\begin{example}\label{ex:g7:negatives:add}
Selangkah demi selangkah, dikelompokkan sesuai kenyamanan:
\begin{align*}
(-8) + (+3) &= -5 && \text{(tanda berlawanan: } 8 - 3 = 5,
\text{ tanda milik } 8)\\
(-2.5) + (-1.5) &= -4 && \text{(tanda sama: jumlahkan jaraknya)}\\
(+6) + (-9) + (+3) &= (-3) + (+3) = 0 .
\end{align*}
\end{example}

\section{Mengurangkan bilangan bertanda}

\begin{theorem}[Aturan pengurangan]\label{thm:g7:negatives:sub}
Mengurangkan sebuah \omterm{def:g7:negatives:def}{bilangan bertanda} sama saja dengan
\emph{menjumlahkan \omterm{def:g7:negatives:def}{lawannya}}:
\[
a - b = a + (-b) .
\]
\end{theorem}

\begin{proof}[Mengapa ini berhasil]
$a - b$ adalah bilangan yang jika ditambahkan pada $b$ memberi $a$.
Periksalah bahwa $a + (-b)$ mengerjakan tugas itu:
$\bigl(a + (-b)\bigr) + b = a + \bigl((-b) + b\bigr) = a + 0 = a$.
\end{proof}

\begin{example}\label{ex:g7:negatives:sub}
Setiap pengurangan menjadi penjumlahan, lalu
\cref{met:g7:negatives:add} berlaku:
\begin{align*}
5 - 9 &= 5 + (-9) = -4, \\
(-3) - (+4) &= (-3) + (-4) = -7, \\
(-2) - (-6) &= (-2) + (+6) = +4 .
\end{align*}
Baris terakhirnya yang terkenal: \emph{mengurangkan bilangan negatif
berarti menjumlahkan}. Menghapus utang membuatmu lebih kaya!
\end{example}

\begin{example}[Selisih suhu]\label{ex:g7:negatives:temperature}
Dalam satu hari suhunya berubah dari $-6$ derajat menjadi $+9$
derajat. Kenaikannya adalah \omterm{ex:g1:subtraction:difference}{selisih}
\[
9 - (-6) = 9 + 6 = 15 \text{ derajat}.
\]
Pada garis bilangan: dari $-6$ naik ke $0$ ada $6$ derajat, dari $0$
ke $9$ ada $9$ derajat lagi.
\end{example}

\section{Koordinat pada seluruh bidang}

\begin{definition}[Koordinat bertanda]\label{def:g7:negatives:coords}
Dengan dua sumbu berskala yang saling \omterm{def:g6:lines:perp}{tegak lurus} dan berpotongan di
titik asal $O$, setiap titik pada bidang memperoleh dua koordinat
bertanda $(x, y)$: $x$ yang negatif berarti di kiri sumbu tegak, $y$
yang negatif berarti di bawah sumbu mendatar.
\end{definition}

\begin{omfigure}
\begin{tikzpicture}[scale=0.62]
  \draw[gray!40, thin] (-4.4,-3.4) grid (4.4,3.4);
  \draw[-Stealth] (-4.6,0) -- (4.6,0) node[below, font=\small] {$x$};
  \draw[-Stealth] (0,-3.6) -- (0,3.6) node[left, font=\small] {$y$};
  \node[below left, font=\small] at (0,0) {$O$};
  \foreach \x in {-4,-2,2,4} \node[below, font=\footnotesize] at (\x,0)
    {$\x$};
  \foreach \y in {-2,2} \node[left, font=\footnotesize] at (0,\y) {$\y$};
  \fill[omDef] (3,2) circle (2.4pt)
    node[above right, font=\small] {$A(3, 2)$};
  \fill[omThm] (-2,3) circle (2.4pt)
    node[above left, font=\small] {$B(-2, 3)$};
  \fill[omMeth] (-3,-2) circle (2.4pt)
    node[below left, font=\small] {$C(-3, -2)$};
  \fill[omProp] (2,-3) circle (2.4pt)
    node[below right, font=\small] {$D(2, -3)$};
\end{tikzpicture}

{\small Satu titik pada setiap perempat bidangnya. Bacalah selalu
koordinat mendatarnya lebih dulu.}
\end{omfigure}

\section{Latihan}

\begin{exercise}[$\star$]\label{exo:g7:negatives:1}
Sebutkan \omterm{def:g7:negatives:def}{lawan} dari: $7$; \ $-3.5$; \ $0$; \ $+12$. Lalu lengkapi:
``dua bilangan saling berlawanan jika \omterm{def:g1:addition:def}{jumlahnya}~?''.
\end{exercise}

\begin{exercise}[$\star$]\label{exo:g7:negatives:2}
Salin dan lengkapi dengan $<$ atau $>$:
\[
-4 \;?\; 1, \qquad
-4 \;?\; -7, \qquad
-2.5 \;?\; -2.4, \qquad
0 \;?\; -0.1 .
\]
\end{exercise}

\begin{exercise}[$\star$]\label{exo:g7:negatives:3}
Urutkan dari yang terkecil ke yang terbesar: $-3$; \ $2.5$; \ $-0.5$; \ $-3.1$;
\ $0$; \ $1$.
\end{exercise}

\begin{exercise}[$\star$]\label{exo:g7:negatives:4}
Hitunglah:
\[
(-5) + (-9), \qquad
(-12) + (+7), \qquad
(+3.5) + (-1.5), \qquad
(-6) + (+6).
\]
\end{exercise}

\begin{exercise}[$\star$]\label{exo:g7:negatives:5}
Ubahlah menjadi penjumlahan, lalu hitunglah:
\[
4 - 11, \qquad
(-5) - (+3), \qquad
(-7) - (-10), \qquad
0 - (-8).
\]
\end{exercise}

\begin{exercise}[$\star$]\label{exo:g7:negatives:6}
Hitunglah dengan mengelompokkan secara cerdik (\omterm{def:g7:negatives:def}{lawan} lebih dulu):
\[
(+7) + (-12) + (+3) + (-8), \qquad
(-5.5) + (+9) + (-4.5) + (+1).
\]
\end{exercise}

\begin{exercise}[$\star$]\label{exo:g7:negatives:7}
Suhunya $-3$ derajat pada pagi ini; suhunya naik $8$ derajat pada
siang hari, lalu turun $12$ derajat pada malam hari. Tulislah satu
perhitungan yang memberi suhu akhirnya, lalu hitunglah.
\end{exercise}

\begin{exercise}[$\star$]\label{exo:g7:negatives:8}
Letakkan pada sistem koordinat: $A(2, 3)$; $B(-3, 1)$; $C(-2, -2)$;
$D(3, -1)$; $E(0, -3)$. Titik mana yang terletak pada sebuah sumbu?
\end{exercise}

\begin{exercise}[$\star\star$]\label{exo:g7:negatives:9}
Hitunglah \omterm{ex:g1:subtraction:difference}{selisih} antara suhu tertinggi dan terendah: Moskwa pada
bulan Januari, tertinggi $-4$ derajat, terendah $-13$ derajat;
Helsinki, tertinggi $2$, terendah $-9$.
\end{exercise}

\begin{exercise}[$\star\star$]\label{exo:g7:negatives:10}
Seorang penyelam berada pada ketinggian $-12$ m (dua belas meter di
bawah permukaan). Ia naik $7$ m, lalu turun $4$ m. Pada ketinggian
berapa ia sekarang? Berapa lagi ia harus naik untuk mencapai permukaan?
\end{exercise}

\begin{exercise}[$\star\star$]\label{exo:g7:negatives:11}
Letakkan $A(-2, 1)$, $B(2, 1)$, $C(2, -2)$. Carilah koordinat titik
$D$ agar $ABCD$ menjadi persegi panjang, lalu hitunglah \omterm{def:g6:measure:perimeter}{kelilingnya}
(bilanglah petaknya untuk panjang sisinya).
\end{exercise}

\begin{exercise}[$\star\star\star$]\label{exo:g7:negatives:12}
Pada persegi ajaib di bawah ini, setiap baris, kolom dan diagonal
harus punya \omterm{def:g1:addition:def}{jumlah} yang sama. Salin dan lengkapilah:
\begin{center}
\renewcommand{\arraystretch}{1.4}
\begin{tabular}{|c|c|c|}
\hline
$?$ & $-7$ & $0$ \\
\hline
$-3$ & $-2$ & $?$ \\
\hline
$-4$ & $?$ & $?$ \\
\hline
\end{tabular}
\end{center}
(Mulailah dengan mencari \omterm{def:g1:addition:def}{jumlah} ajaibnya lewat diagonal yang hampir
lengkap.)
\end{exercise}

\section{Soal: Tahun nol yang tidak pernah ada}

\begin{problem}[{Soal akhir pekan --- bilangan bertanda pada garis
waktu sejarah: lama waktu yang melintasi batas era, trik para
astronom, dan asas perubahan bersih}]
\label{pb:g7:negatives:1}
Garis bilangan pada bab ini punya kembaran yang terkenal di dunia
nyata: garis waktu sejarah, dengan tahun SM memanjang ke kiri dan
tahun M memanjang ke kanan. Tetapi garis milik para sejarawan itu
menyembunyikan perangkap yang merusak banyak perhitungan:
\emph{tidak ada tahun \omterm{def:g1:counting:zero}{nol}} --- tahun 1 SM langsung disusul tahun 1 M.
Soal ini menghitung dengan suhu, lift dan rekening bank, lalu
memperbaiki \omterm{def:g1:counting:zero}{nol} sejarah yang rusak itu dengan cara para astronom:
dengan \omterm{def:g7:negatives:def}{bilangan bertanda}.

\medskip
\noindent\textbf{Bagian I --- Pemanasan, dan perangkap yang pertama.}
\begin{enumerate}
  \item Puncak Mont Blanc berada pada ketinggian $+4\,808$ m; tepi
        Laut Mati pada $-430$ m. Hitunglah \omterm{ex:g1:subtraction:difference}{selisih} ketinggian keduanya.
  \item Sebuah rekening bank memuat $-45$ euro (saldo minus).
        Pemiliknya menyetor $120$ euro, membayar tagihan $60$ euro,
        lalu menyetor $15$ euro. Tulislah saldonya sebagai satu
        rangkaian penjumlahan lalu hitunglah. Setoran tunggal berapa
        yang akan mengembalikan $-45$ semula tepat ke $0$?
  \item Julius Caesar lahir pada $100$ SM dan wafat pada $44$ SM.
        Dengan membilang tahun SM sebagai bilangan negatif ($100$ SM
        sebagai $-100$, $44$ SM sebagai $-44$), hitunglah
        $(-44) - (-100)$: berapa tahun Caesar hidup?
  \item Kaisar Augustus lahir pada $63$ SM dan wafat pada tahun $14$
        M. Resep yang sama memberi $14 - (-63) = 77$ tahun --- tetapi
        para sejarawan bersikeras ia wafat pada umur $76$. Keberatan
        mereka: hitungan $77$ itu berjalan melewati tahun yang tidak
        pernah ada. Tahun yang mana?
  \item Para astronom memperbaiki garis waktunya dengan menamai
        ulang: tahun M tetap dengan bilangannya ($+14$ tetap $+14$),
        tetapi $1$ SM menjadi $0$, $2$ SM menjadi $-1$, dan secara
        umum tahun $n$ SM menjadi $-(n - 1)$. Ubahlah $63$ SM, lalu
        ulangi perhitungan pertanyaan 4 dengan bilangan para astronom
        itu. Apakah jawabannya kini memuaskan para sejarawan?
\end{enumerate}

\medskip
\noindent\textbf{Bagian II --- Menghitung melintasi era.}
Mulai sekarang, pakailah penamaan ulang para astronom untuk setiap
perhitungan, lalu ubah kembali ke SM/M pada akhirnya (tahun yang
berlabel $0$ atau negatif $-m$ adalah tahun $m + 1$ SM milik sejarawan).
\begin{enumerate}[resume]
  \item Ubahlah menjadi bilangan para astronom: $1$ SM; $10$ SM;
        tahun $476$ M; dan $753$ SM, yaitu tahun berdirinya Roma
        menurut legenda.
  \item Berapa tahun berlalu sejak berdirinya Roma ($753$ SM) sampai
        runtuhnya Kekaisaran Barat (tahun $476$ M)?
  \item Sebuah komet kembali setiap $76$ tahun dan muncul pada $12$
        SM. Sebutkan tahun sejarawan untuk tiga kemunculannya
        berikutnya, dan untuk kemunculannya tepat sebelum $12$ SM.
  \item Berapa tahun berlalu dari $5$ SM sampai tahun $5$ M? (Jawaban
        yang menggoda adalah $10$; garis waktunya berkata lain.)
  \item Anak yang lahir pada $3$ SM berumur sepuluh tahun pada tahun
        berapa (SM atau M)?
\end{enumerate}

\medskip
\noindent\textbf{Bagian III --- Asas perubahan bersih.}
\begin{enumerate}[resume]
  \item \emph{Jarak} antara dua bilangan pada garis adalah yang lebih
        besar dikurangi yang lebih kecil --- selalu positif.
        Hitunglah jarak antara $-7$ dan $-2$, lalu antara $-3.5$ dan
        $4.5$.
  \item Sebuah lift bermula dari lantai $0$ lalu melakukan gerakan
        $+3$, $-5$, $+2$, $-1$, $+4$. Hitunglah lantai akhirnya
        dengan mengelompokkan secara cerdik
        (\cref{met:g7:negatives:add}, \omterm{def:g7:negatives:def}{lawan} lebih dulu bila mungkin).
        Apakah melakukan gerakannya dengan urutan yang berbeda akan
        mengubah tujuannya? Mengapa?
  \item Hari demi hari, suhunya berubah $+2$, $-4$, $+1$, $0$, $-3$,
        $+5$, $-2$ derajat selama satu minggu. Hitunglah
        \emph{perubahan bersihnya}. Kalau minggu itu bermula pada
        $-2$ derajat, di mana ia berakhir? Nyatakan asasnya: nilai
        akhir $=$ nilai awal $+$ (\omterm{def:g1:addition:def}{jumlah} semua perubahannya).
  \item Perjalanan seorang pendaki membawanya kembali tepat ke
        ketinggian awalnya. Berapa \omterm{def:g1:addition:def}{jumlah} semua perubahan
        ketinggiannya, dan mengapa? Dari enam perubahan yang
        tercatat, lima di antaranya $+4$, $-7$, $+2$, $+3$, $-5$
        (dalam ratusan meter): carilah yang keenam.
  \item Penutupnya, kembali ke Roma. Seorang sejarawan menghitung:
        ``Republiknya berjalan dari $509$ SM sampai $27$ SM, lalu
        Kekaisarannya sampai tahun $476$ M; seluruhnya, dari $509$ SM
        sampai $476$ M, yaitu $509 + 476 = 985$ tahun negara Romawi.''
        Ulangilah perhitungannya dengan bilangan para astronom. Berapa
        \omterm{def:g1:addition:def}{jumlah} yang benar --- dan tahun yang mana, sekali lagi, yang
        menyebabkan kekeliruannya?

\end{enumerate}
\end{problem}
